%%%
%%% A
%%%
\section*{A}\addcontentsline{toc}{section}{A}\addcontentsline{loh}{figure}{\#\#\#\#\#\#\#\# A}

%%%%%%%%%% 阿 %%%%%%%%%%
\subsection*{阿}\addcontentsline{loh}{figure}{阿 \dpy{a1}}

\begin{EntryWithPhonetic}{阿}{a1}{7}{⾩}
  \definition{pref.}{em dialetos do sul para formar termos carinhosos, antes de nomes de animais de estimação, sobrenomes monossilábicos ou números que denotam ordem de antiguidade em uma família; anexado a 大, 二, 三,\dots\ para indicar a ordem de nascimento, às vezes carrega um sentido de intimidade | antes dos termos de parentesco; na frente de um sobrenome, de um nome próprio ou de um determinado título, com uma conotação de intimidade | em alguns contextos, pode soar infantil ou muito informal, por exemplo, chamar um colega de trabalho por ``阿 + Nome'' sem intimidade}[阿妈===mamãe | 阿明 ===forma carinhosa de chamar uma pessoa cujo nome é Ming]
  \seeref{e1}
\end{EntryWithPhonetic}

\begin{EntryWithPhonetic}{阿非利加洲}{a1fei1li4jia1 zhou1}{7,8,7,5,9}{⾩,⾮,⼑,⼒,⽔}
  \definition{s.}{África; Continente Africano}
\end{EntryWithPhonetic}

\begin{EntryWithPhonetic}{阿哥}{a1ge1}{7,10}{⾩,⼝}
  \definition{s.}{irmão mais velho (afetivo) | forma afetuosa de tratamento entre homens da mesma idade}[阿哥，帮我拿一下书包！===Irmão, ajude"-me com minha mochila escolar!]
\end{EntryWithPhonetic}

\begin{EntryWithPhonetic}{阿克苏}{a1ke4su1}{7,7,7}{⾩,⼗,⾋}
  \definition*{s.}{Aqsu shehiri ou cidade de Aksu | Wilayiti de Aqsu ou prefeitura de Aksu em Xinjiang, 新疆}
  \seealsoref{新疆}{xin1jiang1}
\end{EntryWithPhonetic}

\begin{EntryWithPhonetic}{阿拉伯语}{a1la1bo2yu3}{7,8,7,9}{⾩,⼿,⼈,⾔}[HSK 7-9]
  \definition{s.}{árabe; idioma árabe}
\end{EntryWithPhonetic}

\begin{EntryWithPhonetic}{阿姨}{a1yi2}{7,9}{⾩,⼥}[HSK 4]
  \definition[个,位,名,些]{s.}{tia; uma forma de tratamento para uma mulher da geração dos pais; dirigir"-se a uma mulher que tem aproximadamente a mesma idade da sua mãe, geralmente não é parente | babá em uma família; professora em um jardim de infância | tia; irmã da mãe (mais comum no sul da China)}[阿姨，生日快乐！===Tia, feliz aniversário! | 阿姨，这个苹果多少钱一斤？===Tia/Senhora, quanto custa o quilo dessas maçãs? | 阿姨，我想喝水。===Tia/Babá, eu quero beber água.]
  \synonymref{大娘}{da4niang2}
  \synonymref{大姨}{da4yi2}
  \synonymref{姨妈}{yi2ma1}
  \synonymref{姨娘}{yi2niang2}
  \antonymref{叔叔}{shu1shu5}
  \antonymref{姨父}{yi2fu5}
\end{EntryWithPhonetic}

%%%%%%%%%% 呵 %%%%%%%%%%
\subsection*{呵}\addcontentsline{loh}{figure}{呵 \dpy{a1}}

\begin{EntryWithPhonetic}{呵}{a1}{8}{⼝}
  \variantof{啊}
  \seeref{he1}
\end{EntryWithPhonetic}

%%%%%%%%%% 啊 %%%%%%%%%%
\subsection*{啊}\addcontentsline{loh}{figure}{啊 \dpy{a1}}

\begin{EntryWithPhonetic}{啊}{a1}{10}{⼝}[HSK 4]
  \definition{interj.}{``Ah!''; ``Oh!''; expressa surpresa ou admiração}
  \seeref{a2}
  \seeref{a3}
  \seeref{a4}
  \seeref{a5}
\end{EntryWithPhonetic}

\begin{EntryWithPhonetic}{啊呀}{a1ya1}{10,7}{⼝,⼝}
  \definition{interj.}{``Oh meu Deus!''; interjeição de surpresa}
\end{EntryWithPhonetic}

\begin{EntryWithPhonetic}{啊哟}{a1yo5}{10,9}{⼝,⼝}
  \definition{interj.}{``Meu Deus!''; ``Oh!''; ``Ai!''; interjeição de surpresa ou dor}
\end{EntryWithPhonetic}

\begin{EntryWithPhonetic}{啊}{a2}{10}{⼝}
  \definition{interj.}{``Eh?''; ``Ei?''; ``Que?''; ``Por que?''; indica uma pergunta complementar, uma dúvida ou um pedido de opinião}
  \seeref{a1}
  \seeref{a3}
  \seeref{a4}
  \seeref{a5}
\end{EntryWithPhonetic}

%%%%%%%%%% 嗄 %%%%%%%%%%
\subsection*{嗄}\addcontentsline{loh}{figure}{嗄 \dpy{a2}}

\begin{EntryWithPhonetic}{嗄}{a2}{13}{⼝}
  \variantof{啊}
  \seeref{sha4}
\end{EntryWithPhonetic}

%%%%%%%%%% 啊 %%%%%%%%%%
\subsection*{啊}\addcontentsline{loh}{figure}{啊 \dpy{a3}}

\begin{EntryWithPhonetic}{啊}{a3}{10}{⼝}
  \definition{interj.}{``Eh?''; ``Meu!''; ``E aí?''; ``Que?''; expressa surpresa e dúvida}
  \seeref{a1}
  \seeref{a2}
  \seeref{a4}
  \seeref{a5}
\end{EntryWithPhonetic}

\begin{EntryWithPhonetic}{啊}{a4}{10}{⼝}
  \definition{interj.}{``Bem!''; ``Sim!''; expressa concordância, pronúncia mais curta | ``Oh!''; ``Ah!''; indica que compreendeu, com pronúncia mais longa | ``Oh!''; expressa surpresa ou admiração, com pronúncia mais longa | ``Desgraça!''; expressa tristeza ou pesar}
  \seeref{a1}
  \seeref{a2}
  \seeref{a3}
  \seeref{a5}
\end{EntryWithPhonetic}

\begin{EntryWithPhonetic}{啊}{a5}{10}{⼝}[HSK 2]
  \definition{part.}{usado no final da frase para expressar admiração | usado no final da frase para expressar afirmação, justificativa, insistência, recomendação, etc. | usado no final da frase para indicar dúvida | usado para fazer uma pequena pausa na frase, chamando a atenção para o que vem a seguir | usado após os itens enumerados | usado após verbos repetitivos, indica um processo longo}
  \seeref{a1}
  \seeref{a2}
  \seeref{a3}
  \seeref{a4}
\end{EntryWithPhonetic}

%%%%%%%%%% 哎 %%%%%%%%%%
\subsection*{哎}\addcontentsline{loh}{figure}{哎 \dpy{ai1}}

\begin{EntryWithPhonetic}{哎}{ai1}{8}{⼝}[HSK 7-9]
  \definition{interj.}{``Por que?''; ``Ei!''; ``Ai!''; expressa surpresa ou insatisfação | ``Ei!''; ``Cuidado!''}
\end{EntryWithPhonetic}

\begin{EntryWithPhonetic}{哎呀}{ai1ya1}{8,7}{⼝,⼝}[HSK 7-9]
  \definition{interj.}{expressa surpresa ou espanto | expressa reclamação, impaciência, etc.}
\end{EntryWithPhonetic}

%%%%%%%%%% 哀 %%%%%%%%%%
\subsection*{哀}\addcontentsline{loh}{figure}{哀 \dpy{ai1}}

\begin{EntryWithPhonetic}{哀}{ai1}{9}{⼝}
  \definition*{s.}{Sobrenome: Ai}
  \definition{adj.}{triste; pesaroso}
  \definition{adv.}{tristemente; lamentavelmente}
  \definition{s.}{luto | tristeza; pesar | pena; misericórdia}
  \definition{v.}{lamentar; lamentar"-se por | Literário: estar triste}
  \synonymref{悲}{bei1}
  \antonymref{乐}{le4}
\end{EntryWithPhonetic}

\begin{EntryWithPhonetic}{哀求}{ai1qiu2}{9,7}{⼝,⽔}[HSK 7-9]
  \definition{v.}{suplicar; implorar | suplicar; implorar piedosamente}
  \synonymref{恳求}{ken3qiu2}
  \synonymref{乞求}{qi3qiu2}
  \synonymref{请求}{qing3qiu2}
  \synonymref{要求}{yao1qiu2}
  \antonymref{逼狭}{bi1 xia2}
  \antonymref{逼迫}{bi1po4}
  \antonymref{命令}{ming4ling4}
\end{EntryWithPhonetic}

%%%%%%%%%% 唉 %%%%%%%%%%
\subsection*{唉}\addcontentsline{loh}{figure}{唉 \dpy{ai1}}

\begin{EntryWithPhonetic}{唉}{ai1}{10}{⼝}
  \definition{interj.}{``Sim!''; ``Certo!''; ``Bem!'' | ``Ai de mim!''; o som dos suspiros}
  \seeref{ai4}
\end{EntryWithPhonetic}

%%%%%%%%%% 挨 %%%%%%%%%%
\subsection*{挨}\addcontentsline{loh}{figure}{挨 \dpy{ai1}}

\begin{EntryWithPhonetic}{挨}{ai1}{10}{⼿}
  \definition{prep.}{por turnos; em sequência; indica sequencialmente}
  \definition{v.}{estar próximo de; estar (chegar) perto de; abordar}
  \seeref{ai2}
  \synonymref{依}{yi1}
  \antonymref{打}{da3}
\end{EntryWithPhonetic}

\begin{EntryWithPhonetic}{挨家挨户}{ai1jia1-ai1hu4}{10,10,10,4}{⼿,⼧,⼿,⼾}[HSK 7-9]
  \definition{expr.}{porta a porta; ir de casa em casa; um após o outro}
\end{EntryWithPhonetic}

\begin{EntryWithPhonetic}{挨着}{ai1zhe5}{10,11}{⼿,⽬}[HSK 6]
  \definition{adv.}{ao lado de; perto de; imediatamente depois}
  \synonymref{靠近}{kao4jin4}
  \synonymref{沿着}{yan2zhe5}
  \antonymref{远离}{yuan3li2}
\end{EntryWithPhonetic}

\begin{EntryWithPhonetic}{挨}{ai2}{10}{⼿}[HSK 6]
  \definition{v.}{sofrer; suportar | prolongar; lutar para superar (tempos difíceis); superar (tempo) com dificuldade | atrasar; adiar; procrastinar}
  \seeref{ai1}
  \synonymref{依}{yi1}
  \antonymref{打}{da3}
\end{EntryWithPhonetic}

\begin{EntryWithPhonetic}{挨打}{ai2/da3}{10,5}{⼿,⼿}[HSK 6]
  \definition{v.+compl.}{levar uma surra; ser atacado; ser espancado}
  \antonymref{攻击}{gong1ji1}
\end{EntryWithPhonetic}

%%%%%%%%%% 癌 %%%%%%%%%%
\subsection*{癌}\addcontentsline{loh}{figure}{癌 \dpy{ai2}}

\begin{EntryWithPhonetic}{癌}{ai2}{17}{⽧}[HSK 7-9]
  \definition{s.}{câncer; carcinoma; tumor maligno}
\end{EntryWithPhonetic}

\begin{EntryWithPhonetic}{癌症}{ai2zheng4}{17,10}{⽧,⽧}[HSK 7-9]
  \definition[种]{s.}{câncer; tumores malignos no corpo}
\end{EntryWithPhonetic}

%%%%%%%%%% 矮 %%%%%%%%%%
\subsection*{矮}\addcontentsline{loh}{figure}{矮 \dpy{ai3}}

\begin{EntryWithPhonetic}{矮}{ai3}{13}{⽮}[HSK 4]
  \definition{adj.}{baixo; pequeno | baixo; refere"-se a \emph{status}, posição, classificação, nível, etc.}[他比我矮。===Ele é mais baixo que eu. | 这栋楼很矮，只有三层。===Esse prédio é baixo, tem só três andares. | 她虽然矮，但是跑得很快！===Ela pode ser baixinha, mas corre muito rápido!]
  \seealsoref{低}{di1}
  \synonymref{低}{di1}
  \antonymref{高}{gao1}
\end{EntryWithPhonetic}

\begin{EntryWithPhonetic}{矮凳}{ai3deng4}{13,14}{⽮,⼏}
  \definition{s.}{banquinho baixo | banqueta}[这个矮凳是木制的，很结实。===Este banquinho é de madeira e bem resistente.]
\end{EntryWithPhonetic}

\begin{EntryWithPhonetic}{矮林}{ai3lin2}{13,8}{⽮,⽊}
  \definition{s.}{mata rasteira | bosque baixo}[这片矮林里有很多野兔和鸟类。===Neste bosque baixo há muitos coelhos selvagens e pássaros. | 山坡上长满了矮林，远看像绿色的地毯。===A encosta está coberta de mata rasteira, que de longe parece um tapete verde.]
\end{EntryWithPhonetic}

\begin{EntryWithPhonetic}{矮胖}{ai3pang4}{13,9}{⽮,⾁}
  \definition{adj.}{atarracado; gorducho; rechonchudo; roliço; baixo e robusto | chamar alguém diretamente de 矮胖 pode ser ofensivo}[我家猫矮胖矮胖的，像个毛球。===Meu gato é baixinho e gordinho, parece uma bolinha de pelo.]
  \antonymref{修长}{xiu1chang2}
\end{EntryWithPhonetic}

\begin{EntryWithPhonetic}{矮人}{ai3ren2}{13,2}{⽮,⼈}
  \definition{s.}{anão; pessoa de baixa estatura (indivíduo) | homúnculo; figuras criadas artificialmente pelos alquimistas em frascos de destilação | nanismo}[他虽然是矮人，但很有力气。===Embora ele seja baixo, é muito forte. | 北欧神话中的矮人是技艺高超的工匠。===Na mitologia nórdica, os anões são artesãos habilidosos. | 他因为身高被嘲笑为‘矮人’，这让他很伤心。===Ele foi zombado por ser chamado de ‘anão’ devido à sua altura, o que o magoou.]
\end{EntryWithPhonetic}

\begin{EntryWithPhonetic}{矮树}{ai3shu4}{13,9}{⽮,⽊}
  \definition[棵]{s.}{arbusto | árvore pequena, baixa}[矮树比高树更容易修剪。===Árvores baixas são mais fáceis de podar do que árvores altas. | 我们种了些矮树作为花园的边界。===Plantamos alguns arbustos como cerca natural do jardim.]
\end{EntryWithPhonetic}

\begin{EntryWithPhonetic}{矮小}{ai3xiao3}{13,3}{⽮,⼩}[HSK 4]
  \definition{adj.}{subdimensionado; curto e pequeno; baixo e pequeno | quando usado para pessoas, pode soar depreciativo se não for em contexto neutro ou afetuoso}[这位矮小的老人是村里的智者。===Este idoso baixinho é o sábio da vila. | 这种矮小的灌木适合盆栽。===Este tipo de arbusto pequeno é ideal para vasos. | 山脚下有一片矮小的房屋，显得格外宁静。===Ao pé da montanha, havia casas baixas que transmitiam uma tranquilidade única.]
  \antonymref{高大}{gao1da4}
\end{EntryWithPhonetic}

\begin{EntryWithPhonetic}{矮星}{ai3xing1}{13,9}{⽮,⽇}
  \definition{s.}{estrela anã}[白矮星是恒星演化的最终阶段之一。===Anãs brancas são um dos estágios finais da evolução estelar.]
\end{EntryWithPhonetic}

\begin{EntryWithPhonetic}{矮子}{ai3zi5}{13,3}{⽮,⼦}
  \definition{s.}{pessoa baixa; anão; baixinho}[白雪公主和七个小矮子住在森林里。===Branca de Neve e os sete anões vivem na floresta. | 用``矮子''称呼他人是不礼貌的。===Chamar alguém de ``baixinho'' é falta de educação.]
  \antonymref{巨人}{ju4ren2}
\end{EntryWithPhonetic}

%%%%%%%%%% 艾 %%%%%%%%%%
\subsection*{艾}\addcontentsline{loh}{figure}{艾 \dpy{ai4}}

\begin{EntryWithPhonetic}{艾}{ai4}{5}{⾋}
  \definition*{s.}{Botânica: Artemísia chinesa (Artemisia argyi)}
  \definition*{s.}{Sobrenome: Ai}
  \definition{adj.}{Literário: gracioso; bonito; lindo}
  \definition{s.}{artemísia; absinto; artemísia chinesa}
  \definition{v.}{Literário: parar; terminar}
  \seeref{yi4}
\end{EntryWithPhonetic}

\begin{EntryWithPhonetic}{艾滋病}{ai4zi1bing4}{5,12,10}{⾋,⽔,⽧}[HSK 7-9]
  \definition*{s.}{AIDS, Síndrome da Imunodeficiência Adquirida}
\end{EntryWithPhonetic}

%%%%%%%%%% 唉 %%%%%%%%%%
\subsection*{唉}\addcontentsline{loh}{figure}{唉 \dpy{ai4}}

\begin{EntryWithPhonetic}{唉}{ai4}{10}{⼝}[HSK 7-9]
  \definition{interj.}{``Oh!''; ``Ah!''; ``Bem!''; interjeição que expressa tristeza ou arrependimento | ``Bem!''; ``Argh!''; usado para responder ou reconhecer}
  \seeref{ai1}
\end{EntryWithPhonetic}

%%%%%%%%%% 爱 %%%%%%%%%%
\subsection*{爱}\addcontentsline{loh}{figure}{爱 \dpy{ai4}}

\begin{EntryWithPhonetic}{爱}{ai4}{10}{⽖}[HSK 1]
  \definition*{s.}{Sobrenome: Ai}
  \definition[个]{s.}{amor; afeição profunda; preocupação profunda; especialmente amor entre pessoas}[爱是理解和包容。===O amor é compreensão e tolerância.]
  \definition{v.}{amar; ter sentimentos profundos por pessoas ou coisas | gostar; gostar de; estar interessado em |  cuidar; valorizar; ter em alta estima; cuidar bem de | estar apto a; ter o hábito de}[他们深深爱着对方。===Eles se amam profundamente. | 我爱我的家人。===Eu amo minha família. | 我爱旅行。===Eu adoro viajar.]
  \seealsoref{热爱}{re4'ai4}
  \synonymref{怜}{lian2}
  \synonymref{恋}{lian4}
  \synonymref{热爱}{re4'ai4}
  \synonymref{喜}{xi3}
  \synonymref{喜爱}{xi3'ai4}
  \synonymref{喜欢}{xi3huan5}
  \antonymref{恶}{e4}
  \antonymref{恨}{hen4}
\end{EntryWithPhonetic}

\begin{EntryWithPhonetic}{爱爱}{ai4'ai5}{10,10}{⽖,⽖}
  \definition{v.}{Coloquial: fazer amor ou relações íntimas | pode ser usado como um apelido entre casais, transmitindo ternura | pode soar vulgar se usado em contextos inadequados}[他们俩刚结婚，天天都想爱爱。===Eles acabaram de se casar e querem fazer amor todo dia. | 爱爱，你今天好漂亮！===Amor, você está linda hoje!]
\end{EntryWithPhonetic}

\begin{EntryWithPhonetic}{爱不释手}{ai4bu2shi4shou3}{10,4,12,4}{⽖,⼀,⾤,⼿}[HSK 7-9]
  \definition{expr.}{``Não consigo parar de ler.''; ``Amo tanto que não consigo deixar passar.''; gostar (amar) algo tanto que não se pode suportar separar"-se dele}
\end{EntryWithPhonetic}

\begin{EntryWithPhonetic}{爱戴}{ai4dai4}{10,17}{⽖,⼽}
  \definition{v.}{reverenciar; adorar; amar profundamente e respeitar (a líderes, celebridades, etc.)}
  \seealsoref{敬爱}{jing4'ai4}
  \synonymref{爱护}{ai4hu4}
  \synonymref{保护}{bao3hu4}
  \synonymref{好感}{hao3gan3}
  \synonymref{热爱}{re4'ai4}
  \synonymref{羡慕}{xian4mu4}
  \synonymref{珍惜}{zhen1xi1}
  \synonymref{尊敬}{zun1jing4}
  \synonymref{尊重}{zun1zhong4}
  \antonymref{轻蔑}{qing1mie4}
\end{EntryWithPhonetic}

\begin{EntryWithPhonetic}{爱抚}{ai4fu3}{10,7}{⽖,⼿}
  \definition{s.}{carinho; carícia}
  \definition{v.}{acariciar; afagar; cuidar (com ternura)}[他轻轻爱抚她的头发。===Ele afagou suavemente o cabelo dela. | 母亲爱抚婴儿的脸颊。===A mãe acaricia a bochecha do bebê. | 她爱抚着小猫的耳朵。===Ela acariciou as orelhas do gatinho.]
  \synonymref{爱护}{ai4hu4}
  \synonymref{抚摸}{fu3mo1}
  \antonymref{虐待}{nve4dai4}
\end{EntryWithPhonetic}

\begin{EntryWithPhonetic}{爱国}{ai4guo2}{10,8}{⽖,⼞}[HSK 4]
  \definition{adj.}{patriótico; patriotismo}[爱国是每个公民的责任。===O patriotismo é o dever de todo cidadão. | 这部电影讲述了英雄的爱国故事。===Este filme conta a história patriótica de um herói.]
  \definition{v.}{ser patriota; amar o seu país}
\end{EntryWithPhonetic}

\begin{EntryWithPhonetic}{爱好}{ai4hao4}{10,6}{⽖,⼥}[HSK 1]
  \definition[个,种]{s.}{passatempo; interesse; \emph{hobby}; sentimentos de interesse especial ou afeição por algo | 爱好 é mais usado para atividades regulares (esportes, música), enquanto 喜欢 é para preferências gerais}[他的爱好是收集邮票。===Seu hobby era colecionar selos.  | 我的爱好是读书和旅行。===Meus hobbies são ler e viajar.]
  \definition{v.}{estar interessado em; ter prazer em; ter um forte interesse em algo; ter sentimentos profundos por alguém ou algo}
  \seealsoref{喜欢}{xi3huan5}
  \seealsoref{兴趣}{xing4qu4}
  \synonymref{嗜好}{shi4hao4}
  \synonymref{喜爱}{xi3'ai4}
  \synonymref{喜好}{xi3hao4}
  \synonymref{喜欢}{xi3huan5}
  \synonymref{兴趣}{xing4qu4}
  \antonymref{讨厌}{tao3/yan4}
  \antonymref{厌烦}{yan4fan2}
\end{EntryWithPhonetic}

\begin{EntryWithPhonetic}{爱好者}{ai4hao4zhe3}{10,6,8}{⽖,⼥,⽼}
  \definition{s.}{hobbista; amador; entusiasta; fã; amante (de arte, esportes, etc.)}[他是一位摄影爱好者。===Ele é um entusiasta de fotografia. | 她是位潜水爱好者，经常去东南亚潜水。===Ela é uma mergulhadora amadora e frequentemente mergulha no Sudeste Asiático.  | 我们为书法爱好者创建了一个微信群。===Criamos um grupo no WeChat para amantes de caligrafia.]
\end{EntryWithPhonetic}

\begin{EntryWithPhonetic}{爱护}{ai4hu4}{10,7}{⽖,⼿}[HSK 4]
  \definition{v.}{acalentar; valorizar; salvaguardar; cuidar bem de}[全社会都应爱护老年人。===Toda a sociedade deve tratar os idosos com cuidado e respeito. | 请爱护公园里的小动物。===Por favor, tratem os animais do parque com cuidado.]
  \seealsoref{爱惜}{ai4xi1}
  \seealsoref{保护}{bao3hu4}
  \seealsoref{维护}{wei2hu4}
  \synonymref{爱戴}{ai4dai4}
  \synonymref{爱惜}{ai4xi1}
  \synonymref{保护}{bao3hu4}
  \synonymref{保养}{bao3yang3}
  \synonymref{体贴}{ti3tie1}
  \synonymref{维护}{wei2hu4}
  \synonymref{珍贵}{zhen1gui4}
  \synonymref{珍惜}{zhen1xi1}
  \synonymref{珍重}{zhen1zhong4}
  \synonymref{尊敬}{zun1jing4}
  \antonymref{虐待}{nve4dai4}
  \antonymref{破坏}{po4huai4}
  \antonymref{欺负}{qi1fu5}
  \antonymref{伤害}{shang1hai4}
  \antonymref{损害}{sun3hai4}
  \antonymref{损坏}{sun3huai4}
\end{EntryWithPhonetic}

\begin{EntryWithPhonetic}{爱理不理}{ai4li3-bu4li3}{10,11,4,11}{⽖,⽟,⼀,⽟}[HSK 7-9]
  \definition{expr.}{frio e indiferente; distante}
\end{EntryWithPhonetic}

\begin{EntryWithPhonetic}{爱面子}{ai4/mian4zi5}{10,9,3}{⽖,⾯,⼦}[HSK 7-9]
  \definition{v.+compl.}{estar preocupado em salvar a face; ser vigilante em relação à reputação; ser sensível ao próprio orgulho; valorizar minha própria dignidade e ter medo que os outros me desprezem}[他爱面子，怕别人笑话他。===Ele se importa com sua reputação e tem medo que os outros riam dele.]
\end{EntryWithPhonetic}

\begin{EntryWithPhonetic}{爱情}{ai4qing2}{10,11}{⽖,⼼}[HSK 2]
  \definition[段,种,份]{s.}{amor (entre pessoas); afeição}[爱情是盲目的。===O amor é cego. | 爱情如同玫瑰，美丽却带刺。===O amor é como uma rosa, bela mas com espinhos.  | 这首歌讲述了破碎的爱情故事。===Esta música conta uma história de amor fracassado.]
  \synonymref{恋爱}{lian4'ai4}
\end{EntryWithPhonetic}

\begin{EntryWithPhonetic}{爱人}{ai4ren5}{10,2}{⽖,⼈}[HSK 2]
  \definition[对,个]{s.}{amante; \emph{dollbaby}; namorado(a) | marido ou esposa; mais usado em ocasiões formais}[这是我的爱人。===Este é o meu/minha esposo/companheiro. | 她是我一生的爱人。===Ela é o amor da minha vida. | 请携带爱人出席晚宴。===Por favor, traga seu cônjuge para o jantar.]
  \synonymref{伴侣}{ban4lv3}
  \synonymref{配偶}{pei4'ou3}
  \synonymref{情侣}{qing2lv3}
  \antonymref{仇人}{chou2ren2}
\end{EntryWithPhonetic}

\begin{EntryWithPhonetic}{爱上}{ai4shang4}{10,3}{⽖,⼀}
  \definition{v.}{perder o coração por; apaixonar"-se por}[他在旅行时爱上了一位法国女孩。===Ele se apaixonou por uma garota francesa durante a viagem.  | 我从来没想过自己会爱上健身。===Eu nunca imaginei que iria me apaixonar por academia.]
\end{EntryWithPhonetic}

\begin{EntryWithPhonetic}{爱惜}{ai4xi1}{10,11}{⽖,⼼}[HSK 7-9]
  \definition{v.}{valorizar; prezar; estimar; usar com moderação; não desperdiçar}
  \seealsoref{爱护}{ai4hu4}
  \seealsoref{珍惜}{zhen1xi1}
  \synonymref{爱护}{ai4hu4}
  \synonymref{保护}{bao3hu4}
  \synonymref{敬爱}{jing4'ai4}
  \synonymref{敬重}{jing4zhong4}
  \synonymref{珍贵}{zhen1gui4}
  \synonymref{珍惜}{zhen1xi1}
  \synonymref{珍重}{zhen1zhong4}
  \antonymref{浪费}{lang4fei4}
  \antonymref{殴打}{ou1da3}
\end{EntryWithPhonetic}

\begin{EntryWithPhonetic}{爱心}{ai4xin1}{10,4}{⽖,⼼}[HSK 3]
  \definition[片]{s.}{amor; carinho; compaixão; um sentimento de preocupação e carinho por outras pessoas ou animais}
  \synonymref{关爱}{guan1'ai4}
  \synonymref{关心}{guan1xin1}
  \synonymref{热心}{re4xin1}
  \antonymref{残忍}{can2ren3}
  \antonymref{冷漠}{leng3mo4}
  \antonymref{无情}{wu2qing2}
\end{EntryWithPhonetic}

%%%%%%%%%% 碍 %%%%%%%%%%
\subsection*{碍}\addcontentsline{loh}{figure}{碍 \dpy{ai4}}

\begin{EntryWithPhonetic}{碍}{ai4}{13}{⽯}
  \definition{v.}{atrapalhar; dificultar; obstruir; estar no caminho de | levar em consideração}
\end{EntryWithPhonetic}

\begin{EntryWithPhonetic}{碍事}{ai4/shi4}{13,8}{⽯,⼅}[HSK 7-9]
  \definition{s.}{importa; tem consequências | (usualmente em frases negativas) sem consequência, não importa}[这决定不碍什么事。===Esta decisão não importa.]
  \definition{v.+compl.}{ser um obstáculo; estar no caminho; manter"-se sob os pés de alguém; afetar o trabalho; causar inconveniência}
  \seealsoref{影响}{ying3xiang3}
  \synonymref{妨碍}{fang2'ai4}
  \antonymref{方便}{fang1bian4}
  \antonymref{顺利}{shun4li4}
\end{EntryWithPhonetic}

%%%%%%%%%% 厂 %%%%%%%%%%
\subsection*{厂}\addcontentsline{loh}{figure}{厂 \dpy{an1}}

\begin{EntryWithPhonetic}{厂}{an1}{2}{⼚}[Kangxi 27]
  \definition{s.}{usado principalmente em nomes pessoais}[他名中有个厂字。===O nome dele contém a palavra `An'.]
  \seeref{chang3}
  \seeref{han3}
\end{EntryWithPhonetic}

%%%%%%%%%% 广 %%%%%%%%%%
\subsection*{广}\addcontentsline{loh}{figure}{广 \dpy{an1}}

\begin{EntryWithPhonetic}{广}{an1}{3}{⼴}[Kangxi 53]
  \definition{s.}{mais comum em nomes de pessoas; o mesmo que 庵}[广安是我的朋友。===An'an é meu amigo.]
  \seeref{guang3}
  \seeref{yan3}
  \seealsoref{庵}{an1}
\end{EntryWithPhonetic}

%%%%%%%%%% 安 %%%%%%%%%%
\subsection*{安}\addcontentsline{loh}{figure}{安 \dpy{an1}}

\begin{EntryWithPhonetic}{安}{an1}{6}{⼧}[HSK 4]
  \definition*{s.}{Sobrenome: An}
  \definition{adj.}{pacífico; quieto; tranquilo; calmo | seguro; protegido | com boa saúde | em paz; bem}
  \definition{adv.}{pacificamente; silenciosamente | com segurança; em segurança | em perguntas retóricas: ``como?''}
  \definition{pron.}{usado como pronome interrogativo, como em 哪里，怎么; 谁，何，如何}
  \definition{s.}{segurança; proteção; paz | ampère; abreviação de ampère, 安培}
  \definition{v.}{tranquilizar (a mente de alguém); acalmar | contentar"-se; ficar satisfeito | colocar em uma posição adequada; encontrar um lugar para | instalar; consertar; encaixar; configurar | trazer (uma acusação contra alguém); dar (a alguém um apelido); reivindicar (crédito por algo) | abrigar (uma intenção) | acalmar; estabilizar | sentir"-se satisfeito e à vontade}
  \seealsoref{安培}{an1pei2}
  \seealsoref{何}{he2}
  \seealsoref{哪里}{na3li5}
  \seealsoref{如何}{ru2he2}
  \seealsoref{谁}{shei2}
  \seealsoref{怎么}{zen3me5}
  \antonymref{乱}{luan4}
  \antonymref{危}{wei1}
\end{EntryWithPhonetic}

\begin{EntryWithPhonetic}{安定}{an1ding4}{6,8}{⼧,⼧}[HSK 7-9]
  \definition{adj.}{estável; tranquilo; estabelecido; pacífico e em ordem; estável e normal, sem flutuações}
  \definition{s.}{tranquilizante; medicina ocidental comumente usada, com efeitos sedativos e anticonvulsivantes}
  \definition{v.}{acalmar; estabilizar; manter}
  \seealsoref{安宁}{an1ning2}
  \seealsoref{稳定}{wen3ding4}
  \synonymref{安宁}{an1ning2}
  \synonymref{稳定}{wen3ding4}
  \antonymref{动荡}{dong4dang4}
  \antonymref{扰乱}{rao3luan4}
  \antonymref{骚乱}{sao1luan4}
\end{EntryWithPhonetic}

\begin{EntryWithPhonetic}{安顿}{an1dun4}{6,10}{⼧,⾴}
  \definition{adj.}{tranquilo; pacífico; estável; seguro}
  \definition{v.}{ajudar a se estabelecer; organizar"-se; encontrar um lugar para; tomar as providências adequadas (para pessoas ou coisas) a fim de estabilizá"-las}
  \synonymref{安放}{an1fang4}
  \synonymref{安排}{an1pai2}
  \synonymref{安置}{an1zhi4}
  \synonymref{部署}{bu4shu3}
  \synonymref{布置}{bu4zhi4}
  \synonymref{放置}{fang4zhi4}
  \synonymref{计划}{ji4hua4}
\end{EntryWithPhonetic}

\begin{EntryWithPhonetic}{安放}{an1fang4}{6,8}{⼧,⽅}
  \definition{v.}{colocar em um lugar apropriado; dispor; pôr | colocar em determinado lugar; entronizar | guarde em local seguro}
  \seealsoref{安置}{an1zhi4}
  \synonymref{安顿}{an1dun4}
  \synonymref{安排}{an1pai2}
  \synonymref{安置}{an1zhi4}
  \synonymref{摆放}{bai3fang4}
  \synonymref{部署}{bu4shu3}
  \synonymref{放置}{fang4zhi4}
  \synonymref{计划}{ji4hua4}
  \antonymref{移动}{yi2dong4}
\end{EntryWithPhonetic}

\begin{EntryWithPhonetic}{安抚}{an1fu3}{6,7}{⼧,⼿}[HSK 7-9]
  \definition{v.}{pacificar; consolar; apaziguar; tranquilizar e acalmar}
  \synonymref{安慰}{an1wei4}
  \synonymref{慰问}{wei4wen4}
  \synonymref{欣慰}{xin1wei4}
  \antonymref{催促}{cui1cu4}
  \antonymref{鼓动}{gu3dong4}
  \antonymref{恐吓}{kong3he4}
  \antonymref{威胁}{wei1xie2}
\end{EntryWithPhonetic}

\begin{EntryWithPhonetic}{安徽}{an1hui1}{6,17}{⼧,⼻}
  \definition*{s.}{Anhui, abreviado como , é uma região administrativa de nível provincial; seu nome é uma combinação dos primeiros caracteres das prefeituras de Anqing () e Huizhou () da época}
  \seealsoref{安庆}{an1qing4}
  \seealsoref{徽州}{hui1zhou1}
  \seealsoref{皖}{wan3}
\end{EntryWithPhonetic}

\begin{EntryWithPhonetic}{安家}{an1/jia1}{6,10}{⼧,⼧}
  \definition{v.+compl.}{estabelecer"-se; fixar residência}
\end{EntryWithPhonetic}

\begin{EntryWithPhonetic}{安检}{an1jian3}{6,11}{⼧,⽊}[HSK 6]
  \definition{s.}{verificação de segurança}
  \definition{v.}{realizar verificação de segurança}
\end{EntryWithPhonetic}

\begin{EntryWithPhonetic}{安静}{an1jing4}{6,14}{⼧,⾭}[HSK 2]
  \definition{adj.}{silencioso; tranquilo; sem som; sem barulho e sem algazarra}
  \seealsoref{寂静}{ji4jing4}
  \seealsoref{平静}{ping2jing4}
  \seealsoref{清静}{qing1jing4}
  \synonymref{寂静}{ji4jing4}
  \synonymref{宁静}{ning2jing4}
  \synonymref{平静}{ping2jing4}
  \synonymref{清静}{qing1jing4}
  \antonymref{吵闹}{chao3nao4}
  \antonymref{沸腾}{fei4teng2}
  \antonymref{呼啸}{hu1xiao4}
  \antonymref{热闹}{re4nao5}
  \antonymref{喧哗}{xuan1hua2}
  \antonymref{喧闹}{xuan1nao4}
\end{EntryWithPhonetic}

\begin{EntryWithPhonetic}{安眠药}{an1mian2yao4}{6,10,9}{⼧,⽬,⾋}[HSK 7-9]
  \definition[片,颗,粒,瓶]{s.}{comprimido para dormir; soporífero; pílula para dormir; medicamentos que podem suprimir o córtex cerebral e induzir o sono}
\end{EntryWithPhonetic}

\begin{EntryWithPhonetic}{安宁}{an1ning2}{6,5}{⼧,⼧}[HSK 7-9]
  \definition{adj.}{pacífico; tranquilo; descreve um estado de ordem normal sem fatores externos que causem desordem ou inquietação | calmo; composto; livre de preocupações, ansiedades ou inquietações}
  \seealsoref{安定}{an1ding4}
  \synonymref{安静}{an1jing4}
  \synonymref{安全}{an1quan2}
  \synonymref{和平}{he2ping2}
  \synonymref{冷静}{leng3jing4}
  \synonymref{宁静}{ning2jing4}
  \synonymref{平安}{ping2'an1}
  \synonymref{平静}{ping2jing4}
  \synonymref{舒适}{shu1shi4}
  \synonymref{稳定}{wen3ding4}
  \synonymref{悠闲}{you1xian2}
  \antonymref{动荡}{dong4dang4}
  \antonymref{烦闷}{fan2men4}
  \antonymref{烦恼}{fan2nao3}
  \antonymref{烦躁}{fan2zao4}
  \antonymref{混乱}{hun4luan4}
  \antonymref{骚扰}{sao1rao3}
\end{EntryWithPhonetic}

\begin{EntryWithPhonetic}{安排}{an1pai2}{6,11}{⼧,⼿}[HSK 3]
  \definition{s.}{plano; programação; organização; tabela do plano de atividades ou horários}
  \definition{v.}{organizar (assuntos) de acordo com a sequência ou regras; tratar as coisas de acordo com uma determinada ordem ou regras | atribuir tarefas a alguém; colocar as pessoas nos cargos de trabalho determinados, conforme planejado}
  \seealsoref{安置}{an1zhi4}
  \synonymref{安顿}{an1dun4}
  \synonymref{安置}{an1zhi4}
  \synonymref{部署}{bu4shu3}
  \synonymref{布置}{bu4zhi4}
  \synonymref{打算}{da3suan5}
  \synonymref{计划}{ji4hua4}
  \synonymref{设计}{she4ji4}
  \synonymref{调整}{tiao2zheng3}
  \synonymref{支配}{zhi1pei4}
  \antonymref{混乱}{hun4luan4}
  \antonymref{随意}{sui2/yi4}
\end{EntryWithPhonetic}

\begin{EntryWithPhonetic}{安培}{an1pei2}{6,11}{⼧,⼟}
  \definition{clas.}{A; empréstimo linguístico: ampere; física: unidade de corrente elétrica}
\end{EntryWithPhonetic}

\begin{EntryWithPhonetic}{安庆}{an1qing4}{6,6}{⼧,⼴}
  \definition*{s.}{Anqing, cidade de nível de prefeitura em Anhui (安徽)}
\end{EntryWithPhonetic}

\begin{EntryWithPhonetic}{安全}{an1quan2}{6,6}{⼧,⼊}[HSK 2]
  \definition{adj.}{seguro; protegido; sem perigo; sem ameaças; sem acidentes}
  \definition{s.}{segurança; proteção; refere"-se a um estado ou conceito, geralmente indicando ausência de ameaças ou perigo}
  \seealsoref{风险}{feng1xian3}
  \synonymref{安定}{an1ding4}
  \synonymref{安静}{an1jing4}
  \synonymref{安宁}{an1ning2}
  \synonymref{和平}{he2ping2}
  \synonymref{平安}{ping2'an1}
  \synonymref{平和}{ping2he2}
  \antonymref{危机}{wei1ji1}
  \antonymref{危急}{wei1ji2}
  \antonymref{危险}{wei1xian3}
\end{EntryWithPhonetic}

\begin{EntryWithPhonetic}{安神}{an1/shen2}{6,9}{⼧,⽰}
  \definition{v.+compl.}{acalmar (ou tranquilizar) os nervos; aliviar a inquietação da mente e tranquilizar o corpo | Medicina Chinesa: aliviar o mal"-estar do corpo e da mente}
\end{EntryWithPhonetic}

\begin{EntryWithPhonetic}{安慰}{an1wei4}{6,15}{⼧,⼼}[HSK 5]
  \definition{adj.}{confortar; tranquilizar; consolar; apaziguar}
  \definition[句,通,番,声,个]{s.}{conforto; consolo; comportamento que alivia a dor de alguém e o acalma com palavras ou gestos}
  \definition{v.}{confortar; consolar; acalmar e confortar; deixar a mente tranquila}
  \synonymref{安抚}{an1fu3}
  \synonymref{慰劳}{wei4lao2}
  \synonymref{慰问}{wei4wen4}
  \synonymref{问候}{wen4hou4}
  \synonymref{欣慰}{xin1wei4}
  \antonymref{刺激}{ci4ji1}
  \antonymref{打击}{da3ji1}
  \antonymref{恐吓}{kong3he4}
  \antonymref{威胁}{wei1xie2}
  \antonymref{吓唬}{xia4hu5}
  \antonymref{责备}{ze2bei4}
  \antonymref{指责}{zhi3ze2}
\end{EntryWithPhonetic}

\begin{EntryWithPhonetic}{安稳}{an1wen3}{6,14}{⼧,⽲}[HSK 7-9]
  \definition{adj.}{seguro; suave e estável; estável | composto; calmo e equilibrado; (comportamento) estável; calmo}
  \synonymref{安定}{an1ding4}
  \synonymref{安宁}{an1ning2}
  \synonymref{从容}{cong2rong2}
  \synonymref{巩固}{gong3gu4}
  \synonymref{坚固}{jian1gu4}
  \synonymref{牢固}{lao2gu4}
  \synonymref{稳定}{wen3ding4}
  \synonymref{稳重}{wen3zhong4}
  \synonymref{自在}{zi4zai5}
  \antonymref{动荡}{dong4dang4}
  \antonymref{动摇}{dong4yao2}
  \antonymref{焦躁}{jiao1zao4}
  \antonymref{危急}{wei1ji2}
  \antonymref{摇晃}{yao2huang4}
\end{EntryWithPhonetic}

\begin{EntryWithPhonetic}{安心}{an1xin1}{6,4}{⼧,⼼}[HSK 7-9]
  \definition{adj.}{aliviado; à vontade; tranquilo; seguro}
  \definition{v.}{abrigar (más) intenções; acalentar certas intenções; ter (pensamentos ruins) em mente}
  \seealsoref{放心}{fang4/xin1}
  \seealsoref{心安}{xin1'an1}
  \synonymref{放心}{fang4/xin1}
  \synonymref{坦然}{tan3ran2}
  \antonymref{不安}{bu4'an1}
  \antonymref{操心}{cao1/xin1}
  \antonymref{担心}{dan1/xin1}
  \antonymref{挂念}{gua4nian4}
  \antonymref{苦恼}{ku3nao3}
  \antonymref{忧虑}{you1lv4}
\end{EntryWithPhonetic}

\begin{EntryWithPhonetic}{安逸}{an1yi4}{6,11}{⼧,⾡}[HSK 7-9]
  \definition{adj.}{fácil; fácil e confortável; relaxado e confortável}
  \definition{s.}{conforto e facilidade; conforto}
  \synonymref{安定}{an1ding4}
  \synonymref{安静}{an1jing4}
  \synonymref{安宁}{an1ning2}
  \synonymref{舒畅}{shu1chang4}
  \synonymref{舒适}{shu1shi4}
  \synonymref{舒服}{shu1fu5}
  \synonymref{痛快}{tong4kuai5}
  \synonymref{悠闲}{you1xian2}
  \antonymref{奔波}{ben1bo1}
  \antonymref{艰苦}{jian1ku3}
  \antonymref{艰难}{jian1nan2}
  \antonymref{艰辛}{jian1xin1}
  \antonymref{疲倦}{pi2juan4}
  \antonymref{疲劳}{pi2lao2}
\end{EntryWithPhonetic}

\begin{EntryWithPhonetic}{安置}{an1zhi4}{6,13}{⼧,⽹}[HSK 4]
  \definition{v.}{providenciar; encontrar um lugar para; ajudar a estabelecer"-se; colocar pessoas ou coisas em uma determinada posição ou organizá"-las adequadamente}
  \seealsoref{安放}{an1fang4}
  \seealsoref{安排}{an1pai2}
  \synonymref{安顿}{an1dun4}
  \synonymref{安放}{an1fang4}
  \synonymref{安排}{an1pai2}
  \synonymref{安装}{an1zhuang1}
  \synonymref{部署}{bu4shu3}
  \synonymref{布置}{bu4zhi4}
  \synonymref{放置}{fang4zhi4}
  \synonymref{计划}{ji4hua4}
  \synonymref{睡眠}{shui4mian2}
  \antonymref{丢弃}{diu1qi4}
  \antonymref{抛弃}{pao1qi4}
\end{EntryWithPhonetic}

\begin{EntryWithPhonetic}{安装}{an1zhuang1}{6,12}{⼧,⾐}[HSK 3]
  \definition{v.}{instalar; consertar; configurar; fixar máquinas ou equipamentos (geralmente conjuntos) em um determinado local, de acordo com métodos e especificações específicos}
  \synonymref{安置}{an1zhi4}
  \synonymref{设置}{she4zhi4}
  \synonymref{装配}{zhuang1pei4}
  \synonymref{装置}{zhuang1zhi4}
  \synonymref{组装}{zu3zhuang1}
  \antonymref{拆除}{chai1chu2}
\end{EntryWithPhonetic}

%%%%%%%%%% 庵 %%%%%%%%%%
\subsection*{庵}\addcontentsline{loh}{figure}{庵 \dpy{an1}}

\begin{EntryWithPhonetic}{庵}{an1}{11}{⼴}
  \definition*{s.}{Sobrenome: An}
  \definition[个,座]{s.}{cabana | convento de freiras; templos budistas, principalmente onde vivem as freiras}
\end{EntryWithPhonetic}

%%%%%%%%%% 岸 %%%%%%%%%%
\subsection*{岸}\addcontentsline{loh}{figure}{岸 \dpy{an4}}

\begin{EntryWithPhonetic}{岸}{an4}{8}{⼭}[HSK 5]
  \definition{adj.}{arrogante; orgulhoso; grandioso (de maneira sombria ou condescendente)}
  \definition[条,道,段,面]{s.}{margem; costa; litoral; terreno à beira da água}
\end{EntryWithPhonetic}

\begin{EntryWithPhonetic}{岸上}{an4shang4}{8,3}{⼭,⼀}[HSK 5]
  \definition{s.}{em terra; costa; margem | na margem do rio; na beira do rio}
\end{EntryWithPhonetic}

%%%%%%%%%% 按 %%%%%%%%%%
\subsection*{按}\addcontentsline{loh}{figure}{按 \dpy{an4}}

\begin{EntryWithPhonetic}{按}{an4}{9}{⼿}[HSK 3]
  \definition{prep.}{de acordo com; à luz de; com base em; em conformidade com}
  \definition{v.}{pressionar; empurrar para baixo; pressionar ou apertar com a mão ou os dedos | pôr de parte; deixar de lado; deixar para mais tarde | restringir; controlar; inibir; parar | verificar; consultar | comentar ou anotar (por um editor ou autor)}
  \synonymref{按照}{an4zhao4}
  \synonymref{压}{ya1}
  \synonymref{照}{zhao4}
\end{EntryWithPhonetic}

\begin{EntryWithPhonetic}{按键}{an4jian4}{9,13}{⼿,⾦}[HSK 7-9]
  \definition[个]{s.}{tecla; botão; teclas pressionadas manualmente}[键盘上的按键非常灵敏。===As teclas do teclado são muito responsivas.]
\end{EntryWithPhonetic}

\begin{EntryWithPhonetic}{按理}{an4li3}{9,11}{⼿,⽟}
  \definition{adv.}{de acordo com o princípio ou a razão; no curso normal dos eventos; normalmente | de acordo com a razão; de acordo com a prática comum; por (bons) direitos}
\end{EntryWithPhonetic}

\begin{EntryWithPhonetic}{按理说}{an4li3shuo1}{9,11,9}{⼿,⽟,⾔}[HSK 7-9]
  \definition{adv.}{de acordo com o princípio ou a razão; no curso normal dos eventos; normalmente | é razoável dizer que\dots}
\end{EntryWithPhonetic}

\begin{EntryWithPhonetic}{按摩}{an4mo2}{9,15}{⼿,⼿}[HSK 5]
  \definition{s.}{massagem; empurrar, pressionar, beliscar e amassar o corpo de uma pessoa com as mãos para promover a circulação sanguínea, aumentar a resistência da pele e regular a função dos nervos}
\end{EntryWithPhonetic}

\begin{EntryWithPhonetic}{按期}{an4qi1}{9,12}{⼿,⽉}
  \definition{adv.}{no prazo; na hora certa; dentro do prazo}
  \seealsoref{按时}{an4shi2}
  \synonymref{按时}{an4shi2}
  \synonymref{定期}{ding4qi1}
  \synonymref{准时}{zhun3shi2}
  \antonymref{逾期}{yu2/qi1}
\end{EntryWithPhonetic}

\begin{EntryWithPhonetic}{按时}{an4shi2}{9,7}{⼿,⽇}[HSK 4]
  \definition{adv.}{na hora; no horário; pontualmente; de acordo com o tempo estipulado}
  \seealsoref{按期}{an4qi1}
  \seealsoref{及时}{ji2shi2}
  \synonymref{按期}{an4qi1}
  \synonymref{定时}{ding4shi2}
  \synonymref{准时}{zhun3shi2}
\end{EntryWithPhonetic}

\begin{EntryWithPhonetic}{按说}{an4shuo1}{9,9}{⼿,⾔}[HSK 7-9]
  \definition{adv.}{no curso normal dos eventos; ordinariamente; normalmente | de acordo com o fato (senso comum); refere"-se a falar de acordo com fatos ou senso comum; como uma questão de razão; expressões semelhantes incluem 按理 e 按理说}
  \seealsoref{按理}{an4li3}
  \seealsoref{按理说}{an4li3shuo1}
\end{EntryWithPhonetic}

\begin{EntryWithPhonetic}{按照}{an4zhao4}{9,13}{⼿,⽕}[HSK 3]
  \definition{prep.}{de acordo com; em conformidade com; à luz de; com base em; apresentar os fundamentos ou critérios de julgamento para fazer algo}
  \seealsoref{按}{an4}
  \seealsoref{依照}{yi1zhao4}
  \synonymref{按}{an4}
  \synonymref{服从}{fu2cong2}
  \synonymref{根据}{gen1ju4}
  \synonymref{依据}{yi1ju4}
  \synonymref{依照}{yi1zhao4}
  \synonymref{遵守}{zun1shou3}
  \synonymref{遵循}{zun1xun2}
  \synonymref{遵照}{zun1zhao4}
  \antonymref{随意}{sui2/yi4}
  \antonymref{违背}{wei2bei4}
  \antonymref{违反}{wei2fan3}
\end{EntryWithPhonetic}

%%%%%%%%%% 案 %%%%%%%%%%
\subsection*{案}\addcontentsline{loh}{figure}{案 \dpy{an4}}

\begin{EntryWithPhonetic}{案}{an4}{10}{⽊}
  \definition{s.}{mesa; escrivaninha; mesa longa | caso; caso de direito (legal) | registro; arquivo; arquivo de caso | um plano submetido para consideração; proposta; um documento que propõe planos, sugestões, métodos, etc.}
\end{EntryWithPhonetic}

\begin{EntryWithPhonetic}{案件}{an4jian4}{10,6}{⽊,⼈}[HSK 7-9]
  \definition[个,起,件,类]{s.}{caso; caso de direito; caso legal; contencioso e eventos ilegais}
\end{EntryWithPhonetic}

%%%%%%%%%% 暗 %%%%%%%%%%
\subsection*{暗}\addcontentsline{loh}{figure}{暗 \dpy{an4}}

\begin{EntryWithPhonetic}{暗}{an4}{13}{⽇}[HSK 4]
  \definition{adj.}{escuro; opaco; sem graça; pouca luz | escondido; secreto; não revelado | pouco claro; nebuloso; vago; confuso | subterrâneo}
  \definition{adv.}{secretamente | no escuro}
  \seealsoref{黑}{hei1}
  \synonymref{黑}{hei1}
  \synonymref{昏}{hun1}
  \synonymref{幽}{you1}
  \antonymref{亮}{liang4}
  \antonymref{明}{ming2}
\end{EntryWithPhonetic}

\begin{EntryWithPhonetic}{暗地里}{an4di4li3}{13,6,7}{⽇,⼟,⾥}[HSK 7-9]
  \definition{adv.}{secretamente; interiormente; às escondidas}
\end{EntryWithPhonetic}

\begin{EntryWithPhonetic}{暗恋}{an4lian4}{13,10}{⽇,⼼}
  \definition{s.}{amor secreto}
  \definition{v.}{estar secretamente apaixonado por}
  \synonymref{迷恋}{mi2lian4}
  \synonymref{喜欢}{xi3huan5}
\end{EntryWithPhonetic}

\begin{EntryWithPhonetic}{暗杀}{an4sha1}{13,6}{⽇,⽊}[HSK 7-9]
  \definition{s.}{assassinato}
  \definition{v.}{assassinar | matar secretamente}
  \synonymref{谋害}{mou2hai4}
\end{EntryWithPhonetic}

\begin{EntryWithPhonetic}{暗示}{an4shi4}{13,5}{⽇,⽰}[HSK 4]
  \definition[种]{s.}{sugestão; insinuação; intimação; Psicologia: refere"-se ao uso de palavras, gestos, expressões, etc. para fazer as pessoas aceitarem involuntariamente uma determinada opinião ou fazerem algo}
  \definition{v.}{dar uma dica; sugerir secretamente; indicar algo a alguém usando outras palavras, expressões faciais ou gestos sem dizer em voz alta}
\end{EntryWithPhonetic}

\begin{EntryWithPhonetic}{暗香}{an4xiang1}{13,9}{⽇,⾹}
  \definition{s.}{fragrância sutil}
\end{EntryWithPhonetic}

\begin{EntryWithPhonetic}{暗中}{an4zhong1}{13,4}{⽇,⼁}[HSK 7-9]
  \definition{adv.}{no escuro; na escuridão | em segredo; às escondidas; sorrateiramente | furtivamente; nos bastidores}
  \synonymref{黑暗}{hei1'an4}
  \antonymref{公开}{gong1kai1}
\end{EntryWithPhonetic}

\begin{EntryWithPhonetic}{暗自}{an4zi4}{13,6}{⽇,⾃}
  \definition{adv.}{interiormente; secretamente; para si mesmo}
\end{EntryWithPhonetic}

%%%%%%%%%% 昂 %%%%%%%%%%
\subsection*{昂}\addcontentsline{loh}{figure}{昂 \dpy{ang2}}

\begin{EntryWithPhonetic}{昂}{ang2}{8}{⽇}
  \definition{adj.}{alto; subindo}
  \definition{v.}{manter (a cabeça) erguida | elevar; levantar; olhar para cima}
\end{EntryWithPhonetic}

\begin{EntryWithPhonetic}{昂贵}{ang2gui4}{8,9}{⽇,⾙}[HSK 7-9]
  \definition{adj.}{caro; dispendioso; algo é muito caro, o preço é particularmente alto; metaforicamente, o custo de fazer algo é particularmente alto}
  \seealsoref{贵}{gui4}
  \synonymref{高价}{gao1jia4}
  \synonymref{贵重}{gui4zhong4}
  \antonymref{廉价}{lian2jia4}
  \antonymref{便宜}{pian2yi5}
\end{EntryWithPhonetic}

%%%%%%%%%% 凹 %%%%%%%%%%
\subsection*{凹}\addcontentsline{loh}{figure}{凹 \dpy{ao1}}

\begin{EntryWithPhonetic}{凹}{ao1}{5}{⼐}[HSK 7-9]
  \definition{adj.}{afundado; amassado | côncavo; oco; amassado; mais baixo que a área circundante}
  \definition{v.}{cavar; chanfrar | desabar; afundar}
  \seeref{wa1}
  \antonymref{凸}{tu1}
\end{EntryWithPhonetic}

%%%%%%%%%% 熬 %%%%%%%%%%
\subsection*{熬}\addcontentsline{loh}{figure}{熬 \dpy{ao1}}

\begin{EntryWithPhonetic}{熬}{ao1}{14}{⽕}
  \definition{v.}{ensopar; ferver; cozinhar em água}
  \seeref{ao2}
  \synonymref{挨}{ai2}
  \synonymref{炖}{dun4}
  \synonymref{忍}{ren3}
  \synonymref{煮}{zhu3}
  \antonymref{享受}{xiang3shou4}
\end{EntryWithPhonetic}

\begin{EntryWithPhonetic}{熬}{ao2}{14}{⽕}[HSK 7-9]
  \definition{v.}{ferver; ensopar; fazer uma decocção; cozinhar em fogo baixo por muito tempo | preparar; infundir; extrair a essência por fervura longa | resistir; suportar (angústia, tempos difíceis, etc.)}
  \seeref{ao1}
\end{EntryWithPhonetic}

\begin{EntryWithPhonetic}{熬夜}{ao2/ye4}{14,8}{⽕,⼣}[HSK 7-9]
  \definition{v.+compl.}{ficar acordado a noite toda ou até tarde da noite}
  \synonymref{通宵}{tong1xiao1}
\end{EntryWithPhonetic}

%%%%%%%%%% 傲 %%%%%%%%%%
\subsection*{傲}\addcontentsline{loh}{figure}{傲 \dpy{ao4}}

\begin{EntryWithPhonetic}{傲}{ao4}{12}{⼈}[HSK 7-9]
  \definition{adj.}{orgulhoso; altivo | arrogante}
  \definition{v.}{recusar"-se a ceder; desafiar}
\end{EntryWithPhonetic}

\begin{EntryWithPhonetic}{傲慢}{ao4man4}{12,14}{⼈,⼼}[HSK 7-9]
  \definition{adj.}{altivo; arrogante; autoritário}
  \synonymref{高傲}{gao1'ao4}
  \antonymref{谦虚}{qian1xu1}
  \antonymref{谦逊}{qian1xun4}
  \antonymref{虚心}{xu1xin1}
\end{EntryWithPhonetic}

%%%%%%%%%% 奥 %%%%%%%%%%
\subsection*{奥}\addcontentsline{loh}{figure}{奥 \dpy{ao4}}

\begin{EntryWithPhonetic}{奥}{ao4}{12}{⼤}
  \definition*{s.}{Oersted, a unidade eletromagnética de intensidade do campo magnético; abreviação de 奥斯特 | Sobrenome: Ao}
  \definition{adj.}{profundo e difícil de entender; abstruso | significado profundo, não é fácil de entender}
  \definition{s.}{canto secreto da casa; antigamente, referia"-se ao canto sudoeste de uma casa e também, de modo geral, à profundidade de uma casa}
  \seealsoref{奥斯特}{ao4si1te4}
\end{EntryWithPhonetic}

\begin{EntryWithPhonetic}{奥林匹克运动会}{ao4lin2pi3ke4yun4dong4hui4}{12,8,4,7,7,6,6}{⼤,⽊,⼖,⼗,⾡,⼒,⼈}
  \definition*{s.}{Jogos Olímpicos, Olimpíadas}
\end{EntryWithPhonetic}

\begin{EntryWithPhonetic}{奥秘}{ao4mi4}{12,10}{⼤,⽲}[HSK 7-9]
  \definition[个]{s.}{enigma; mistério profundo; fenômenos ou princípios profundos e misteriosos}
  \seealsoref{秘密}{mi4mi4}
  \synonymref{机密}{ji1mi4}
  \synonymref{秘密}{mi4mi4}
  \synonymref{奇妙}{qi2miao4}
  \synonymref{神秘}{shen2mi4}
  \synonymref{微妙}{wei1miao4}
  \antonymref{常识}{chang2shi2}
\end{EntryWithPhonetic}

\begin{EntryWithPhonetic}{奥斯特}{ao4si1te4}{12,12,10}{⼤,⽄,⽜}
  \definition{s.}{Oersted}
\end{EntryWithPhonetic}

\begin{EntryWithPhonetic}{奥特曼}{ao4te4man4}{12,10,11}{⼤,⽜,⽈}
  \definition*{s.}{Ultraman,  super"-herói de ficção científica japonesa}
\end{EntryWithPhonetic}

\begin{EntryWithPhonetic}{奥运}{ao4yun4}{12,7}{⼤,⾡}
  \definition*[届,次]{s.}{Jogos Olímpicos, Olimpíadas; abreviação de 奥林匹克运动会}
  \seealsoref{奥林匹克运动会}{ao4lin2pi3ke4yun4dong4hui4}
\end{EntryWithPhonetic}

\begin{EntryWithPhonetic}{奥运会}{ao4yun4hui4}{12,7,6}{⼤,⾡,⼈}[HSK 7-9]
  \definition*[届,次]{s.}{Jogos Olímpicos, Olimpíadas; abreviação de 奥林匹克运动会}
  \seealsoref{奥林匹克运动会}{ao4lin2pi3ke4yun4dong4hui4}
\end{EntryWithPhonetic}

%%%%%%%%%% 澳 %%%%%%%%%%
\subsection*{澳}\addcontentsline{loh}{figure}{澳 \dpy{ao4}}

\begin{EntryWithPhonetic}{澳}{ao4}{15}{⽔}
  \definition*{s.}{Abreviação de Austrália, 澳大利亚 | Sobrenome: Ao}
  \definition{s.}{baía; uma entrada do mar; um lugar curvo na costa onde os barcos podem ser atracados, frequentemente usado em nomes de lugares}
  \seealsoref{澳大利亚}{ao4da4li4ya4}
\end{EntryWithPhonetic}

\begin{EntryWithPhonetic}{澳大利亚}{ao4da4li4ya4}{15,3,7,6}{⽔,⼤,⼑,⼆}
  \definition*{s.}{Austrália}
\end{EntryWithPhonetic}

%%%%% EOF %%%%%

